% !TeX program = pdflatex
%==============================================================================
%  Lista Teórica 05 --- Métodos Não Paramétricos: Aspectos Teóricos
%  O gabarito sai de "Lista teorica 05 - gabarito.tex".
%==============================================================================
\providecommand{\gabaritoopt}{}
\documentclass[11pt,a4paper]{article}
\usepackage[\gabaritoopt]{../../recursos/latex/estilo-lista}

\begin{document}
\cabecalholista{05}{Métodos Não Paramétricos: Aspectos Teóricos}

%------------------------------------------------------------------------------
\begin{exercicio}[o preço de cada dimensão]
O Teorema da Aula~05 diz que o KNN com $k$ bem escolhido atinge risco
\[
  \E\big[(\rhat_k(\X)-r(\X))^2\big] \;\le\; K'\,n^{-2/(2+d)} .
\]
Suponha que, em $d=1$, $n=100$ observações bastem para atingir um risco alvo
$\delta$ (tratando a desigualdade como igualdade).
\begin{enumerate}[label=(\alph*),leftmargin=1.1cm]
  \item Isolando $n$, mostre que $n=(K'/\delta)^{(2+d)/2}$.
  \item Use o dado de $d=1$ para achar $K'/\delta$ e calcule quantas observações
        seriam necessárias, para o \emph{mesmo} risco, em $d=2$, $d=5$, $d=10$ e
        $d=20$.
  \item O crescimento é polinomial ou exponencial em $d$? Justifique olhando a
        fórmula.
  \item Um colega diz: ``mas com $n$ grande o KNN converge para o classificador
        ótimo em qualquer dimensão, isso é um teorema''. Ele está certo. Por que,
        então, o resultado acima é uma má notícia?
\end{enumerate}
\end{exercicio}

\begin{solucao}
\textbf{(a)} De $\delta=K'n^{-2/(2+d)}$ vem $n^{2/(2+d)}=K'/\delta$ e, elevando
a $(2+d)/2$,
\[
  n = \Big(\frac{K'}{\delta}\Big)^{(2+d)/2}.
\]

\smallskip
\textbf{(b)} Em $d=1$ o expoente é $3/2$, então
$100=(K'/\delta)^{3/2}$ e $K'/\delta=100^{2/3}\approx 21{,}54$. Substituindo:
\begin{center}\small
\renewcommand{\arraystretch}{1.3}
\begin{tabular}{@{}lccccc@{}}
\toprule
$d$ & $1$ & $2$ & $5$ & $10$ & $20$ \\
expoente $(2+d)/2$ & $1{,}5$ & $2$ & $3{,}5$ & $6$ & $11$ \\
\midrule
$n$ necessário & $100$ & $464$ & $4{,}6\times10^{4}$ & $\mathbf{10^{8}}$ & $4{,}6\times10^{14}$ \\
\bottomrule
\end{tabular}
\end{center}
Cem observações em $d=1$ viram \textbf{cem milhões} em $d=10$. Em $d=20$ seriam
$4{,}6\times10^{14}$, cerca de sessenta mil vezes a população da Terra.

\smallskip
\textbf{(c)} \textbf{Exponencial.} O $d$ aparece no \emph{expoente}: escrevendo
$c=K'/\delta$, temos $n=c^{(2+d)/2}=c\cdot(\sqrt{c}\,)^{d}$. Cada dimensão nova
multiplica $n$ por um fator constante $\sqrt{c}$ --- aqui, $\sqrt{21{,}54}\approx
4{,}64$. É exatamente o que a tabela mostra: de $d=10$ para $d=20$, dez
multiplicações por $4{,}64$, ou seja um fator de $4{,}6$ milhões.

\smallskip
\textbf{(d)} As duas afirmações são compatíveis porque falam de coisas
diferentes. A consistência é um resultado \emph{assintótico}: garante que o erro
vai a zero quando $n\to\infty$, mas não diz \emph{a que velocidade}. A taxa
$n^{-2/(2+d)}$ diz a velocidade --- e em $d=20$ o expoente é $-1/11$, de modo que
para dividir o erro por 2 é preciso multiplicar $n$ por $2^{11}=2048$.

Em outras palavras: a garantia existe, mas o $n$ que a torna útil não existe.
Consistência é uma condição \emph{necessária} para levar um método a sério, não
uma razão para usá-lo em dimensão alta.
\end{solucao}

%------------------------------------------------------------------------------
\begin{exercicio}[a vizinhança que deixou de ser vizinha]
Espalhe $n$ pontos uniformemente no cubo $[0,1]^d$. Para que uma bola de raio
$\rho$ em torno de um ponto contenha em média $k$ deles, é preciso
$n\rho^d\approx k$, ou seja
\[
  \rho \;\approx\; \Big(\frac{k}{n}\Big)^{1/d} .
\]
\begin{enumerate}[label=(\alph*),leftmargin=1.1cm]
  \item Com $k=10$ e $n=10\,000$, calcule $\rho$ para $d=1$, $2$, $5$, $10$ e
        $20$. (Repare que $k/n=0{,}001$ em todos os casos.)
  \item Interprete o valor de $d=20$ em palavras: que fração da largura do cubo
        o raio cobre? O que resta da palavra ``vizinho''?
  \item Ainda em $d=20$, quantas observações seriam necessárias para que os $10$
        vizinhos mais próximos coubessem num raio de $\rho=0{,}1$?
  \item Ligue este exercício ao Exercício~3 da Lista Teórica~04, sobre o viés de
        fronteira. Que relação há entre os dois fenômenos?
\end{enumerate}
\end{exercicio}

\begin{solucao}
\textbf{(a)} Com $k/n=10^{-3}$, temos $\rho=(10^{-3})^{1/d}=10^{-3/d}$:
\begin{center}\small
\renewcommand{\arraystretch}{1.3}
\begin{tabular}{@{}lccccc@{}}
\toprule
$d$ & $1$ & $2$ & $5$ & $10$ & $20$ \\
\midrule
$\rho$ & $0{,}001$ & $0{,}032$ & $0{,}251$ & $0{,}501$ & $\mathbf{0{,}708}$ \\
\bottomrule
\end{tabular}
\end{center}

\smallskip
\textbf{(b)} Em $d=20$ o raio necessário é $0{,}708$: a bola cobre cerca de
\textbf{71\% da largura do cubo em cada direção}. Um ``vizinho'' a essa distância
está tão longe de $\x$ quanto um ponto qualquer sorteado ao acaso.

E aí o método perde o argumento que o justifica. O KNN estima $r(\x)$ pela média
dos $Y$ dos vizinhos porque supõe que, se $\x_i$ está perto de $\x$, então
$r(\x_i)$ está perto de $r(\x)$ --- é a hipótese de Lipschitz. Com vizinhos a
$71\%$ da largura do domínio, essa hipótese não diz nada de útil, e o viés do
Teorema, que vale $(k/n)^{2/d}$, deixa de ser pequeno.

\smallskip
\textbf{(c)} Queremos $\rho=(10/n)^{1/20}=0{,}1$, isto é
$10/n=0{,}1^{20}=10^{-20}$, logo
\[
  n = 10\times 10^{20} = \mathbf{10^{21}} .
\]
Um sextilhão de observações --- mais do que grãos de areia na Terra --- só para
que dez vizinhos caibam num décimo da largura do cubo.

\smallskip
\textbf{(d)} São o mesmo fenômeno visto de dois ângulos. Na Lista~04 vimos que o
viés do Nadaraya--Watson explode nos pontos de fronteira, porque ali a janela
fica assimétrica e todos os vizinhos caem do mesmo lado. Aqui vemos que, em
dimensão alta, a janela necessária é tão grande que \emph{todo} ponto está
efetivamente na fronteira do domínio: com $\rho=0{,}708$ num cubo de lado $1$,
não existe ponto cuja bola de vizinhos caiba inteira dentro do cubo.

Ou seja, o defeito que em $d=1$ é localizado nas duas pontas passa, em $d$
grande, a ser o comportamento típico em todo lugar. É a mesma maldição descrita
duas vezes.
\end{solucao}

%------------------------------------------------------------------------------
\begin{exercicio}[de onde sai o expoente]
A demonstração do Teorema chega a
\[
  \E\big[(\rhat_k(\X)-r(\X))^2\big]
   \;\le\; \underbrace{K\Big(\frac{k}{n}\Big)^{2/d}}_{\text{viés}^2}
        + \underbrace{\frac{\sigma^2}{k}}_{\text{variância}} .
\]
\begin{enumerate}[label=(\alph*),leftmargin=1.1cm]
  \item Explique, em uma frase cada, por que o primeiro termo cresce com $k$ e o
        segundo decresce.
  \item Minimize a soma em $k$ e mostre que o ótimo satisfaz
        $k \asymp n^{2/(2+d)}$. \emph{Sugestão:} derive em $k$ e iguale a zero, ou
        simplesmente iguale os dois termos --- a menos de constantes dá o mesmo.
  \item Substitua de volta e obtenha a taxa $n^{-2/(2+d)}$.
  \item O que acontece com o $k$ ótimo quando $d$ cresce, com $n$ fixo? Comente o
        que isso significa na prática.
\end{enumerate}
\end{exercicio}

\begin{solucao}
\textbf{(a)} O \textbf{viés$^2$} cresce com $k$ porque pedir mais vizinhos obriga
a ir buscá-los mais longe, e quanto mais longe estão os $\x_i$, menos $r(\x_i)$
se parece com $r(\x)$ --- é o fator $(k/n)^{2/d}$, que é o raio ao quadrado. A
\textbf{variância} decresce com $k$ porque $\rhat_k$ é a média de $k$ respostas
condicionalmente independentes, e a variância de uma média de $k$ termos é
$\sigma^2/k$.

\smallskip
\textbf{(b)} Igualando os dois termos, que é o atalho usual (o mínimo de uma soma
de uma função crescente e uma decrescente está, a menos de constantes, onde elas
se cruzam):
\[
  \Big(\frac{k}{n}\Big)^{2/d} \asymp \frac1k
  \;\Longrightarrow\;
  k^{2/d}\,k \asymp n^{2/d}
  \;\Longrightarrow\;
  k^{(2+d)/d} \asymp n^{2/d}
  \;\Longrightarrow\;
  k \asymp n^{2/(2+d)} .
\]
Pela derivada dá o mesmo: $\frac{d}{dk}\big[Kn^{-2/d}k^{2/d}+\sigma^2k^{-1}\big]
= \frac{2K}{d}n^{-2/d}k^{2/d-1}-\sigma^2k^{-2}=0$, o que leva a
$k^{2/d+1}=\frac{d\sigma^2}{2K}n^{2/d}$ e à mesma ordem.

\smallskip
\textbf{(c)} Levando $k\asymp n^{2/(2+d)}$ ao termo de variância (o mais fácil
dos dois):
\[
  \frac{\sigma^2}{k} \asymp \sigma^2\,n^{-2/(2+d)} .
\]
No termo de viés dá o mesmo, por construção --- foi assim que escolhemos $k$.
Logo a soma é da ordem de $n^{-2/(2+d)}$.

\smallskip
\textbf{(d)} O expoente $2/(2+d)$ vai a zero quando $d\to\infty$, então
$k\asymp n^{2/(2+d)}\to n^0=$ constante. Ou seja: em dimensão alta, o $k$ ótimo
\textbf{para de crescer com $n$}.

O que isso significa é desanimador. Em dimensão baixa, ganhar mais dados permite
usar mais vizinhos, e a variância cai. Em dimensão alta os vizinhos adicionais
estão tão longe que o viés que eles trazem anula o ganho de variância --- então o
método fica preso a uma média sobre um punhado de pontos, por mais dados que
você colete. É a tradução, em termos de $k$, da taxa lenta do Exercício~1.
\end{solucao}

%------------------------------------------------------------------------------
\begin{exercicio}[como escapar: aditividade]
A maldição vale para métodos que estimam $r$ \emph{sem nenhuma suposição de
forma}. Se estivermos dispostos a supor que $r$ é \textbf{aditiva},
\[
  r(\x) = f_1(x_1) + f_2(x_2) + \dots + f_d(x_d),
\]
cada $f_j$ é uma função de \textbf{uma} variável e pode ser estimada à taxa
unidimensional $n^{-2/3}$. Somando os $d$ erros, o risco fica da ordem de
$d\,n^{-2/3}$.
\begin{enumerate}[label=(\alph*),leftmargin=1.1cm]
  \item Com $d=10$ e $n=10\,000$, calcule as duas taxas: $n^{-2/(2+d)}$ (sem
        suposição) e $d\,n^{-2/3}$ (aditiva). Qual a razão entre elas?
  \item Quantas observações o método sem suposição precisaria para alcançar o
        erro que o aditivo alcança com $10\,000$?
  \item Onde entra $d$ em cada uma das duas taxas? É essa a diferença essencial.
  \item Qual é o preço da suposição? Dê um exemplo de $r$ que o modelo aditivo
        não consegue representar, por mais dados que se tenha.
\end{enumerate}
\end{exercicio}

\begin{solucao}
\textbf{(a)} Com $n=10^4$ e $d=10$:
\[
  n^{-2/(2+d)} = (10^4)^{-1/6} = 10^{-2/3} \approx \mathbf{0{,}215},
  \qquad
  d\,n^{-2/3} = 10\times(10^4)^{-2/3} = 10\times 10^{-8/3} \approx \mathbf{0{,}0215}.
\]
A razão é exatamente \textbf{10}: a suposição de aditividade divide o erro por
dez, com os mesmos dados e sem mudar nada no experimento.

\smallskip
\textbf{(b)} Queremos $n^{-1/6}=0{,}0215$, ou seja $n=0{,}0215^{-6}=10^{10}$.
\textbf{Dez bilhões} de observações para igualar o que a suposição entrega com
dez mil --- um fator de um milhão.

\smallskip
\textbf{(c)} Na taxa sem suposição, $d$ está no \textbf{expoente}:
$n^{-2/(2+d)}$. É por isso que ela degrada exponencialmente. Na taxa aditiva, $d$
é um \textbf{fator multiplicativo}: $d\cdot n^{-2/3}$, com expoente $-2/3$ fixo,
o mesmo de um problema unidimensional. Dobrar $d$ dobra o erro, em vez de exigir
que $n$ seja elevado a uma potência maior.

Essa é a lição geral do curso, e ela não é sobre aditividade em particular:
\textbf{a maldição da dimensionalidade é uma consequência de não supor nada}.
Toda estrutura que você esteja disposto a supor --- aditividade, linearidade,
esparsidade, suavidade de ordem maior --- compra uma taxa melhor. A pergunta
nunca é ``como escapar da maldição'', é ``que suposição eu aceito fazer''.

\smallskip
\textbf{(d)} O preço são as \textbf{interações}. O modelo aditivo não representa
nenhuma função em que o efeito de $x_1$ dependa do valor de $x_2$. O exemplo
canônico é o produto,
\[
  r(x_1,x_2) = x_1 x_2 ,
\]
que não se escreve como $f_1(x_1)+f_2(x_2)$. Um caso concreto no espírito da
Aula~11: o efeito de uma dose de medicamento sobre a pressão pode ser positivo
num grupo de pacientes e negativo em outro; um modelo aditivo é obrigado a
escolher um único efeito médio para a dose e vai errar nos dois grupos.

Note o contraste com o Exercício~1: lá o problema era o $n$ ser insuficiente,
o que mais dados resolveriam. Aqui o viés é da \emph{classe de modelos}, e não
some com $n$ nenhum.
\end{solucao}


\end{document}
